\documentclass[../../../include/open-logic-section]{subfiles}

\begin{document}

\olfileid{sth}{card-arithmetic}{opps}
\olsection{تعریف عمل‌های پایه}

چون لازم نیست ترتیب را پی بگیریم، تعریف حساب کاردینالی بسیار آسان‌تر از
حساب ترتیبی است. جمع، ضرب و توان‌رسانی را هم‌زمان تعریف می‌کنیم.

\begin{defn}
اگر~$\cardfont{a}$ و~$\cardfont{b}$ کاردینال باشند:
\begin{align*}
	\cardfont{a} \cardplus \cardfont{b} &\defis
	\card{\cardfont{a} \disjointsum \cardfont{b}}\\
	\cardfont{a} \cardtimes \cardfont{b} &\defis
	\card{\cardfont{a} \times \cardfont{b}}\\
	\cardexpo{\cardfont{a}}{\cardfont{b}} &\defis
	\card{\funfromto{\cardfont{b}}{\cardfont{a}}}
\end{align*}
که در آن
$\funfromto{X}{Y}=\Setabs{f}{f\text{ تابعی }X\to Y\text{ است}}$.
(به‌آسانی می‌توان نشان داد~$\funfromto{X}{Y}$ برای هر دو مجموعهٔ~$X$ و
$Y$ وجود دارد؛ اثبات را به‌عنوان تمرین واگذار می‌کنیم.)
\end{defn}

\begin{prob}
در~$\Zminus$ اثبات کنید~$\funfromto{X}{Y}$ برای هر دو مجموعهٔ~$X$ و~$Y$
وجود دارد. در~$\ZF$ و به شیوهٔ
\olref[sth][ord-arithmetic][using-addition]{rankcomputation}،
$\setrank{\funfromto{X}{Y}}$ را بر حسب~$\setrank{X}$ و~$\setrank{Y}$
محاسبه کنید.
\end{prob}

شاید توضیح این تعریف سودمند باشد. دربارهٔ جمع: این تعریف از مفهوم جمع
مجزا، یعنی~$\disjointsum$، که در
\olref[ord-arithmetic][add]{defdissum} تعریف شد استفاده می‌کند؛ و
به‌آسانی می‌توان دید در حالت‌های متناهی پاسخ درست می‌دهد. دربارهٔ ضرب:
\olref[sfr][set][pai]{cardnmprod} می‌گوید اگر~$A$ دارای~$n$ عضو و~$B$
دارای~$m$ عضو باشد، $A\times B$ دارای~$n\cdot m$ عضو است؛ پس تعریف ما
صرفاً این اندیشه را به ضرب ترامتناهی تعمیم می‌دهد. توان‌رسانی نیز مشابه
است: صرفاً اندیشه‌ای متناهی را به حالت ترامتناهی تعمیم می‌دهیم. در واقع،
حساب کاردینالی ترامتناهی از برخی جهات بسیار بیش از حساب ترتیبی ترامتناهی
به حساب «معمولی» شباهت دارد:

\begin{prop}\ollabel{cardplustimescommute}
$\cardplus$ و~$\cardtimes$ جابه‌جایی و شرکت‌پذیرند.
\end{prop}

\begin{proof}
برای جابه‌جایی، بنا بر
\olref[cardinals][cardsasords]{lem:CardinalsBehaveRight} کافی است توجه
کنیم
$\cardeq{(\cardfont{a}\disjointsum\cardfont{b})}
{(\cardfont{b}\disjointsum\cardfont{a})}$ و
$\cardeq{(\cardfont{a}\times\cardfont{b})}
{(\cardfont{b}\times\cardfont{a})}$. شرکت‌پذیری را به‌عنوان تمرین واگذار
می‌کنیم.
\end{proof}

\begin{prob}
اثبات کنید~$\cardplus$ و~$\cardtimes$ شرکت‌پذیرند.
\end{prob}

\begin{prop}
$A$ نامتناهی است اگر و تنها اگر
$\card{A}\cardplus1=1\cardplus\card{A}=\card{A}$.
\end{prop}

\begin{proof}
همانند~\olref[cardinals][classing]{generalinfinitycharacter}، این نتیجه
از~\olref[ord-arithmetic][using-addition]{ordinfinitycharacter} و
\olref[cardinals][cardsasords]{lem:CardinalsBehaveRight} حاصل می‌شود.
\end{proof}

این امر توضیح می‌دهد چرا باید برای جمع و ضرب ترتیبی و کاردینالی از
نمادهای متفاوتی استفاده کنیم: این‌ها واقعاً عمل‌هایی \emph{متفاوت‌اند}.
دو نتیجهٔ بعدی نشان می‌دهند توان‌رسانی ترتیبی و کاردینالی نیز عمل‌هایی
متفاوت‌اند. (به یاد آورید~\olref[z][infinity-again]{defnomega} مستلزم
$2=\{0,1\}$ است):

\begin{lem}\ollabel{lem:SizePowerset2Exp}
برای هر~$A$، $\card{\Pow{A}}=\cardexpo{2}{\card{A}}$.
\end{lem}

\begin{proof}
برای هر زیرمجموعهٔ~$B\subseteq A$، بگذارید
$\chi_B\in\funfromto{A}{2}$ چنین تعریف شود:
\begin{align*}
	\chi_{B}(x) &\defis
	\begin{cases}
		1 & \text{اگر }x\in B\\
		0 & \text{در غیر این صورت.}
	\end{cases}
\end{align*}
اکنون بگذارید~$f(B)=\chi_B$؛ این یک~!!a{bijection} به‌صورت
$f\colon\Pow{A}\to\funfromto{A}{2}$ تعریف می‌کند. پس
$\cardeq{\Pow{A}}{\funfromto{A}{2}}$. در نتیجه
$\cardeq{\Pow{A}}{\funfromto{\card{A}}{2}}$، و بنابراین
$\card{\Pow{A}}=\card{\funfromto{\card{A}}{2}}=
\cardexpo{2}{\card{A}}$.
\end{proof}

این اثبات کوتاه اساساً بحث~\olref[sfr][siz][red-alt]{sec} را دربر
می‌گیرد. آنجا نشان دادیم چگونه می‌توان ناشمارابودن~$\Pow{\omega}$ را به
ناشمارابودن مجموعهٔ رشته‌های دودویی نامتناهی، یعنی~$\Bin^\omega$، فرو
کاست. در واقع، $\Bin^\omega$ همان~$\funfromto{\omega}{2}$ است؛ و اثبات
پیشین نشان داد استدلال~\olref[sfr][siz][red-alt]{sec} با هر مجموعهٔ~$A$
به‌جای~$\omega$ نیز کار می‌کند. این نتیجه واقعیت کوتاه زیر را نیز دربارهٔ
توان‌رسانی کاردینالی به دست می‌دهد:

\begin{cor}\ollabel{cantorcor}
برای هر کاردینال~$\cardfont{a}$،
$\cardfont{a}<\cardexpo{2}{\cardfont{a}}$.
\end{cor}

\begin{proof}
از قضیهٔ کانتور (\olref[sfr][siz][car]{thm:cantor}) و
\olref{lem:SizePowerset2Exp}.
\end{proof}
\noindent
پس~$\omega<\cardexpo{2}{\omega}$. اما توجه کنید این نتیجه دربارهٔ
توان‌رسانی \emph{کاردینالی} است. باید آن را با توان‌رسانی
\emph{ترتیبی} مقایسه کرد، زیرا در حالت دوم
$\omega=\ordexpo{2}{\omega}$ (بنگرید به
\olref[ord-arithmetic][expo]{sec}).

حال که دربارهٔ توان‌رسانی کاردینالی سخن می‌گوییم، می‌توانیم «شیوهٔ»
!!{nonenumerable} بودن~$\Real$ را نیز اندکی دقیق‌تر بیان کنیم.

\begin{thm}\ollabel{continuumis2aleph0}
$\card{\Real}=\cardexpo{2}{\omega}$.
\end{thm}

\begin{proof}[طرح اثبات]
راه‌های فراوانی برای اثبات این نتیجه وجود دارد. سرراست‌ترین راه آن است
که~$\cardle{\Pow{\omega}}{\Real}$ و
$\cardle{\Real}{\Pow{\omega}}$ را اثبات کنیم؛ سپس با استفاده از قضیهٔ
شرودر--برنشتاین نتیجه بگیریم
$\cardeq{\Real}{\Pow{\omega}}$، و با
\olref[card-arithmetic][opps]{lem:SizePowerset2Exp} نتیجه بگیریم
$\card{\Real}=\cardexpo{2}{\omega}$. تعریف دو تزریق
$f\colon\Pow{\omega}\to\Real$ و
$g\colon\Real\to\Pow{\omega}$ را به‌عنوان تمرینی روشنگر واگذار می‌کنیم.
\end{proof}

\begin{prob}
اثبات~\olref[sth][card-arithmetic][opps]{continuumis2aleph0} را با
نشان‌دادن~$\cardle{\Pow{\omega}}{\Real}$ و
$\cardle{\Real}{\Pow{\omega}}$ کامل کنید.
\end{prob}

\end{document}
